%!TEX program = pdflatex
\documentclass[11pt, a4paper]{article}

% ─── Packages ────────────────────────────────────────────────────────────────
\usepackage[T1]{fontenc}
\usepackage[utf8]{inputenc}
\usepackage{lmodern}
\usepackage[margin=2.5cm]{geometry}
\usepackage{amsmath, amssymb, bm}
\usepackage{graphicx}
\usepackage{booktabs}
\usepackage{multirow}
\usepackage{array}
\usepackage{xcolor}
\usepackage{hyperref}
\usepackage{natbib}
\usepackage{caption}
\usepackage{subcaption}
\usepackage{float}
\usepackage{algorithm}
\usepackage{algpseudocode}
\usepackage{enumitem}
\usepackage{parskip}
\usepackage{titlesec}
\usepackage{setspace}

\onehalfspacing
\setlength{\parindent}{0pt}
\setlength{\parskip}{6pt}

\hypersetup{
    colorlinks=true,
    linkcolor=blue!60!black,
    citecolor=green!50!black,
    urlcolor=blue!70!black
}

\graphicspath{{figures/}{../exp032c_figures/}{../exp033_figures/}{../exp034_figures/}{../exp035_figures/}{../phase15_figures/}{../phase16d_figures/}}

% ─── Title ───────────────────────────────────────────────────────────────────
\title{
    \vspace{-1cm}
    {\large Capstone Design Project --- Final Report}\\[0.8em]
    \textbf{The Effect of Reward Function Design in RL-Based\\
    BTC Grid Trading: Pareto Frontier, Winner Reversal,\\
    and Out-of-Distribution Robustness of Prospect-Theoretic Reward}
}

\author{
    Chanhee Lee\\
    Department of Computer Science\\
    Seoul National University of Science and Technology\\
    \texttt{happilyfly@seoultech.ac.kr}
}

\date{June 2026}

% ─── Document ────────────────────────────────────────────────────────────────
\begin{document}

\maketitle
\thispagestyle{empty}

% ─── Abstract ────────────────────────────────────────────────────────────────
\begin{abstract}
We quantitatively study the effect of \emph{reward function design} on PPO-based
reinforcement-learning policies for ATR-proportional grid trading on the BTC/USDT
1-hour market. Four reward variants---symmetric (\texttt{sym}), asymmetric
($\beta=3.42$, \texttt{asym}), Differential Sharpe Ratio
($\eta=1/28\,\mathrm{h}$, \texttt{dsr}; \citealt{moody2001dsr}), and
prospect-theoretic ($\alpha=0.68$, $\lambda=3.30$, \texttt{pt};
\citealt{kahneman1979prospect})---are each trained with 10 seeds $\times$ 1M steps
and evaluated on three environments: (i) single-split validation
(Val 2021--2023), (ii) 6-fold combinatorial purged cross-validation
\citep{lopezdeprado2018advances} 15 paths, and (iii) the unsealed
out-of-sample Test set (2024--2026, BTC \$42K\,$\to$\,\$76K bull regime).

We report five key findings.
\textbf{(1)} All four variants exceed the ATR baseline by $+11$--$24\%$ on Val,
yet partition into two clusters \{\texttt{sym}, \texttt{dsr}\} (aggressive) vs.\
\{\texttt{asym}, \texttt{pt}\} (conservative); within-cluster Cohen's $|d|<0.3$
versus across-cluster $|d|>0.8$.
\textbf{(2)} Mechanism: reward loss asymmetry directly determines trade frequency
and holding time (policy distance ratio 2.22$\times$), and the sliding-window
memory structure of DSR results in 2--6$\times$ higher hold rate than other
variants.
\textbf{(3)} Winner reverses across evaluation environments: Val=\texttt{sym}(1.87)
$\to$ CPCV=\texttt{dsr}(1.41, $p<0.001$) $\to$ Test=\texttt{pt}(0.34--0.37,
consistent across two model sources, $p<0.002$).
\textbf{(4)} The OOS robustness of \texttt{pt} arises from a short-holding
mechanism: loss aversion $\lambda=3.30$ combined with concave gain $\alpha=0.68$
trains policies to exit within mean 1.4\,h (max 6\,h), thereby avoiding the
sell-side timing risk of the unseen bull regime. In contrast, \texttt{dsr} learns
long holding (mean 4.58\,h, max 169\,h = 7 days) and records the worst OOS
performance---demonstrating that the same reward formulation underlies both
in-sample advantage and OOS failure as two sides of the same coin.
\textbf{(5)} The Val--Test distribution shift is highly significant
(KS $p<10^{-10}$, $-27\%$ variance), while policies' actions remain nearly
invariant between Val and Test ($\Delta \leq 5\%$), indicating that the
regime shift, not the policy, drives OOS differences.

Our contribution is twofold. First, we demonstrate quantitatively that reward
design exerts its effect through a \emph{trade-off dimension in risk profile}
rather than a single Sharpe ``alpha source.'' Second, we provide the first
quantitative evidence that Kahneman-Tversky prospect theory confers
out-of-distribution safety on RL trading policies. Code and data:
\url{https://github.com/cosmicpotato2047/capstone-rl-trading}.
\end{abstract}

\vspace{1em}
\noindent\textbf{Keywords:} Reinforcement Learning, Grid Trading, PPO,
Reward Design, Prospect Theory, Differential Sharpe Ratio,
Combinatorial Purged CV, Out-of-Distribution Generalization, Cryptocurrency

\newpage
\tableofcontents
\newpage

% ─── 1. Introduction ─────────────────────────────────────────────────────────
\section{Introduction}
\label{sec:intro}

Cryptocurrency markets exhibit micro-oscillations in hourly price that are
largely orthogonal to long-term trends, motivating \emph{grid trading} strategies
that harvest revenue from volatility itself rather than from directional
prediction. Grid trading is a simplified form of market making
\citep{avellaneda2008market}, ranging from fixed-grid rules to volatility-adaptive
variants that scale grid width by the Average True Range (ATR)
\citep{wilder1978atr}. This work studies the latter: an ATR-proportional grid
system whose grid width and limit-order placement are determined by a
reinforcement-learning policy.

\subsection{Motivation and Research Question}
\label{sec:intro_motivation}

An earlier phase of this project hypothesized that PPO RL would discover dynamic
volatility adaptation and outperform a fixed-grid baseline. Post-hoc behavioral
analysis, however, revealed that the trained policy converged to a near-constant
action $[1.0, 0.0]$ (\S\ref{sec:phase1_negative}), and that the
Bayesian-optimized ATR coefficient itself was the principal source of
performance. This negative finding suggested that the symmetric-reward $+$
ATR-proportional combination structurally absorbs the policy's degrees of
freedom, motivating the new research question of this paper:

\begin{quote}
\textbf{In BTC grid trading, how does the design of the reward function affect
the behavioral pattern and generalization performance of RL policies? In
particular, do loss-asymmetric (asymmetric, prospect-theoretic) or
risk-adjusted-return (DSR) reward formulations yield meaningful differences in
single-metric advantage or out-of-sample robustness?}
\end{quote}

\subsection{Contributions}
\label{sec:intro_contribution}

This work makes the following four contributions.

\begin{enumerate}[leftmargin=*]
\item \textbf{A statistically rigorous comparison framework for reward variants.}
Four reward functions (\texttt{sym}, \texttt{asym}, \texttt{dsr}, \texttt{pt}) are
trained in an identical environment with 10 seeds $\times$ 1M steps, and pairwise
statistical tests (Cohen's $d$, bootstrap P(A$>$B), IQM, 5\% CVaR) are applied
following \citet{henderson2018rlreproducibility}'s reproducibility guidelines.
\item \textbf{Scenario D --- Pareto-frontier discovery.} Rather than a single
``best variant'' winner, the four variants statistically partition into two
clusters \{\texttt{sym}, \texttt{dsr}\} vs.\ \{\texttt{asym}, \texttt{pt}\}
forming a Pareto-like frontier on the Sharpe--MDD plane
(\S\ref{sec:pareto}). This outcome falls into none of the pre-registered scenario
branches (A optimistic / B neutral / C pessimistic).
\item \textbf{Causal chain: reward $\to$ behavior $\to$ outcome.} Policy-level
differences across variants are quantified on 1.04M step trajectories via a
policy-distance matrix (within-cluster 0.129 vs.\ across-cluster 0.286), regime-
specific trade/hold statistics, and action distribution comparison, confirming
the mechanism reward loss asymmetry $\to$ reduced trade frequency $\to$
conservative cluster (\S\ref{sec:mechanism}).
\item \textbf{Winner reversal across evaluation environments and OOS robustness
of prospect theory.} The winner reverses completely across single-split Val
(\texttt{sym} 1\textsuperscript{st}) $\to$ CPCV multi-split
(\texttt{dsr} 1\textsuperscript{st}) $\to$ the unsealed Test
(\texttt{pt} 1\textsuperscript{st}, consistent across two sources $p<0.002$)
(\S\ref{sec:robustness}). The OOS robustness of \texttt{pt} is mechanistically
attributed via trajectory analysis to its loss aversion $\lambda=3.30$ and
concave utility $\alpha=0.68$ training policies to exit within mean 1.4\,h,
thereby avoiding the sell-side timing risk of the unseen bull regime. This is,
to our knowledge, the first quantitative demonstration that the Kahneman-Tversky
prospect theory \citep{kahneman1979prospect} confers OOS safety on RL trading
policies.
\end{enumerate}

\subsection{Paper Structure}
\label{sec:intro_structure}

Section \ref{sec:related} surveys related work on RL trading, reward design
theory, and evaluation methodology. Section \ref{sec:method} details the
environment design (MDP, ATR formula, four reward variants) and PPO training.
Section \ref{sec:phase1_negative} briefly reproduces the Phase 1 negative
finding (``RL $=$ Fixed $[1.0, 0.0]$'') as motivation. Section \ref{sec:pareto}
reports the two-cluster Pareto-frontier result on Val (Scenario D), while
Section \ref{sec:mechanism} quantifies the mechanism. Section
\ref{sec:robustness} presents the robustness trio: slippage
(\S\ref{sec:slippage}), CPCV multi-split (\S\ref{sec:cpcv}), and the central
contribution of this paper---the Test OOS evaluation and the mechanistic basis
of prospect theory's robustness (\S\ref{sec:oos}). Section \ref{sec:discussion}
discusses distribution shift quantification and limitations, and Section
\ref{sec:conclusion} concludes.

% ─── 2. Related Work ─────────────────────────────────────────────────────────
\section{Related Work}
\label{sec:related}

This paper sits at the intersection of four areas: (i) market making and grid
trading, (ii) reinforcement-learning-based trading, (iii) reward-design theory,
and (iv) evaluation methodology for trading policies. This section surveys the
key prior work in each area and clarifies where our contribution lies.

\subsection{Market Making and Grid Trading}
\label{sec:related_grid}

The academic ancestor of grid trading is the market-making model of
\citet{avellaneda2008market}. Their analysis derives a closed-form optimal
policy for a risk-averse market maker who faces inventory risk and must place
\emph{asymmetric} bid/ask quotes around the mid-price. Two key insights are
that (a) the bid--ask spread is a function of volatility and remaining time
horizon, and (b) inventory skew shifts the quote midpoint asymmetrically. Grid
trading is a simplified practical variant of this model in which orders are
placed on a pre-defined ladder of price levels in order to monetize
micro-volatility.

The standard mechanism for adapting grid widths to market volatility is the
Average True Range \citep{wilder1978atr} (ATR; 14-period moving average is the
default). ATR is a smoothed average of the true price range and provides a
simple, robust rule that automatically widens the grid in volatility-cluster
regimes. ATR-proportional grids of the form $\Delta = c \cdot \mathrm{ATR}$
(with proportional coefficient $c$) are widely used in practice and serve as
the baseline of this paper.

Most prior academic work in this area is limited to the analysis of
\emph{rule-based} policies, with relatively few attempts to learn the grid-width
decision from data. This paper adopts a structure in which the proportional
coefficient and the number of grid splits in the ATR formula are dynamically
determined by an RL policy, while quantitatively analyzing how the design of
the reward function influences the resulting learned policy's behavior---a
combination, to our knowledge, that has not been explored.

\subsection{Reinforcement Learning for Trading}
\label{sec:related_drl}

The standard DRL baseline for cryptocurrency trading is the PPO/A2C/DQN
comparison framework of \citet{zhang2020drl}, which feeds price time series
directly into a PPO policy and shows consistent improvements over simple
buy-and-hold (B\&H) and moving-average crossover strategies. The ACM survey of
\citet{sun2023rl} and the more recent review of \citet{liu2025review} both
identify two open problems in the area: \emph{(a) interpretability of policy
behavior and (b) generalization across market regimes}. This paper addresses
both directly: \S\ref{sec:mechanism} quantifies the causal chain from reward
form to behavioral difference, and \S\ref{sec:robustness} documents
out-of-sample (OOS) generalization phenomena that are not captured by
single-split validation alone.

There are four directly preceding works. \citet{liu2021bitcoin} feeds the
output of an LSTM price predictor into the state of a PPO trading policy in a
two-stage pipeline. The limitation of this approach is that prediction errors
propagate directly into trading decisions; this paper bypasses the limitation
by \emph{eliminating} the price-prediction stage and using a volatility-adaptive
grid structure instead. \citet{yasin2024rl} compares DQN/PPO/A2C with 20
technical indicators, but the action space is discrete buy/sell with no hold
action. This paper adopts a continuous two-dimensional action (aggressiveness
and profit target) that supports fine-grained grid control.
\citet{pham2025dqn} uses DQN to select one of five predefined strategies as a
meta-strategy, but the strategies themselves are not learned.
\citet{bandarupalli2025risk} augments PPO with a volatility penalty and a
transaction-cost term, but stays at the level of hyperparameter tuning within
a single reward variant and does not compare across reward \emph{forms}, which
is the focus of this paper.

A particularly relevant prior work is \citet{gort2022backtest}, which
empirically shows that \emph{single-split train/test results do not guarantee
OOS generalization} in cryptocurrency DRL trading, and recommends CPCV
(Combinatorial Purged Cross-Validation) and DSR (Deflated Sharpe Ratio). This
paper adopts those recommendations, using Val (2021--2023), CPCV (6-fold inside
the Train partition), and Test (2024-- unsealed) environments concurrently.
Our \emph{winner reversal} finding (Val \texttt{sym} $\to$ CPCV \texttt{dsr}
$\to$ Test \texttt{pt}) is a quantitative confirmation that Gort et al.'s
concern reappears in the dimension of reward-variant selection.

In one sentence, our position is: whereas prior DRL trading literature focuses
on comparing algorithms and state designs under a \emph{single} reward function,
this paper performs a controlled comparison of \emph{four reward functions} on
top of the same algorithm (PPO) and the same environment, and separately
measures their effects on in-sample / multi-split / OOS environments.

\subsection{Reward-Design Theory}
\label{sec:related_reward}

The four reward variants compared in this paper draw from three theoretical
lineages in reinforcement learning and behavioral economics. (i) Symmetric
(\texttt{sym}) is the standard PnL-return reward and serves as our baseline.
(ii) Asymmetric (\texttt{asym}) places a positive weight $\beta>1$ on losses,
a simple first-order approximation of a risk-averse utility. (iii) Differential
Sharpe Ratio (\texttt{dsr}) is the formulation of \citet{moody2001dsr}.
(iv) Prospect-theoretic reward (\texttt{pt}) directly uses the utility function
of \citet{kahneman1979prospect} and \citet{tversky1992cumulative},
$v(x) = \mathrm{sgn}(x) |x|^\alpha$ for $x \geq 0$ and $-\lambda |x|^\alpha$
for $x < 0$, as a per-step reward.

The DSR formulation of \citet{moody2001dsr} uses the \emph{differential change
of the online Sharpe-ratio estimator} as the step reward, tracking first and
second moments via exponentially-weighted moving averages (EWMA, decay rate
$\eta$). This has two implications for RL training: (a) it directly optimizes
a risk-adjusted return rather than raw PnL, and (b) the sliding-window memory
imposes a \emph{multi-step temporal dependency} on the policy. As we show in
\S\ref{sec:mechanism}, this memory structure is what causes the \texttt{dsr}
policy to learn 2--6$\times$ longer holding times than the other variants---a
property that becomes the in-sample advantage of \S\ref{sec:robustness} and
the OOS failure of \S\ref{sec:oos} as two sides of the same causal coin.

Prospect theory \citep{kahneman1979prospect} is a foundational model of
human decision making under risk, with two characteristic parameters:
(a) loss aversion $\lambda \approx 2.25$ (standard value from
\citealt{tversky1992cumulative}) and (b) the concavity exponent
$\alpha \approx 0.88$. While the theory has been widely used to describe
human trader behavior (e.g., the disposition effect of clinging to losses and
realizing gains too early), \emph{direct use of prospect-theoretic utility as
an RL reward} is rare in the literature. To our knowledge, this paper is the
first to directly implement the \citet{kahneman1979prospect} utility as an RL
reward (with $\alpha=0.68$, $\lambda=3.30$, tuned by Optuna) and quantitatively
establish the OOS robustness that the resulting policy exhibits.

\citet{ng1999reward}'s policy-invariance theorem of potential-based reward
shaping states that ``a monotone transformation of the reward signal does not
alter the policy learned by RL.'' At first glance, this paper's comparison of
four reward variants may appear to conflict with the theorem, but it does not:
the \texttt{sym} $\to$ \texttt{asym}/\texttt{pt} transformations are
\emph{non-potential-based} (they apply an asymmetric scale on the negative
half-line), while \texttt{dsr} is not a simple per-step reward transformation
but a \emph{state-extended} reward (the sliding-window moments effectively
become part of the state). Our reward-variant comparison thus lies outside the
scope of the Ng et al.\ theorem, and indeed we confirm in
\S\ref{sec:mechanism} that policy behavior differs across variants in a
statistically significant manner.

\subsection{Evaluation Methodology}
\label{sec:related_eval}

The two principal pitfalls in evaluating RL trading policies are (a) backtest
overfitting and (b) seed dependence of results.
\citet{henderson2018rlreproducibility} empirically demonstrates that seed
variability alone can produce large differences for the same algorithm, and
proposes 5--10 seed statistical comparisons, IQM (Interquartile Mean), and
bootstrap confidence intervals as standard recommendations. This paper uses
10 seeds per variant (Val/exp033) and 100 models (Test), and applies pairwise
Cohen's $d$ and bootstrap $\mathrm{P}(A>B)$ to all comparisons.

\citet{lopezdeprado2018advances} (Ch.\ 8) proposes \emph{Combinatorial Purged
Cross-Validation} (CPCV), which avoids train/test leakage in time-series data
while leveraging $\binom{N}{k}$ split combinations to strengthen the
generalization guarantees of single-split results. We adopt 6-fold CPCV
(15 paths) in \S\ref{sec:cpcv}.

\citet{lopezdeprado2014dsr}'s \emph{Deflated Sharpe Ratio}
(DSR\textsuperscript{*}) is a statistical-significance test on the Sharpe
ratio that corrects for the number of backtest trials and for the skewness/kurtosis
of the underlying distribution. We apply Bonferroni 4-way correction together
with DSR\textsuperscript{*} testing to the multiple comparisons of CPCV
results (\S\ref{sec:cpcv}), and confirm that all four reward variants are
significant relative to the ATR baseline at $p<0.001$.

The methodological distinctiveness of this paper is the \emph{simultaneous use
of three evaluation environments} (Val / CPCV / Test) that allowed us to
discover the winner reversal. Studies relying solely on single-split or solely
on CPCV cannot expose OOS generalization failure, and studies relying solely
on OOS cannot expose the in-sample conservatism of \texttt{pt}. The
simultaneous application of three environments is, in our view, a methodological
contribution: it generalizes the multi-evaluation recommendations of
\citet{henderson2018rlreproducibility} and \citet{gort2022backtest} to the
single dimension of reward-variant comparison.

% ─── 3. Method ───────────────────────────────────────────────────────────────
\section{Method}
\label{sec:method}

\subsection{MDP Formulation}
\label{sec:mdp}

We cast single-asset dynamic grid trading on the BTC/USDT 1-hour series as a
finite-horizon Markov decision process $\mathcal{M} = (\mathcal{S},
\mathcal{A}, P, R, \gamma)$. An episode runs for 720 hours (30 days) starting
from a randomly chosen step within the training/validation range, and each
step $t$ corresponds to a 1-hour price bar.

\paragraph{State $\mathcal{S} \subset \mathbb{R}^7$.}
The observation $\mathbf{s}_t = (s_t^{(0)}, \ldots, s_t^{(6)})$ couples market
and portfolio state:

\begin{align}
s_t^{(0)} &= \log\!\left(p_t \,\big/\, \overline{p}_{t, 168}\right) &&
\text{(price vs.\ 1-week mean)} \\
s_t^{(1)} &= \frac{\bar{p}^{\mathrm{avg}}_t - p_t}{\bar{p}^{\mathrm{avg}}_t} &&
\text{(divergence from average cost)} \\
s_t^{(2)} &= \frac{h_t \cdot p_t}{C_0} &&
\text{(holdings value / capital)} \\
s_t^{(3)} &= \frac{c_t}{C_0} &&
\text{(cash ratio)} \\
s_t^{(4)} &= \frac{\mathrm{ATR}_{168}(p)}{p_t} &&
\text{(volatility)} \\
s_t^{(5)} &= \frac{p_t - p_{t-72}}{p_{t-72}} &&
\text{(short-term trend, 3d)} \\
s_t^{(6)} &= \frac{p_t - p_{t-720}}{p_{t-720}} &&
\text{(long-term trend, 30d)}
\end{align}

Here $p_t$ is the bar close, $\bar{p}^{\mathrm{avg}}_t$ is the average cost of
the current position (set to 0 when no position is held, in which case
$s_t^{(1)}$ is also 0), $h_t$ is the BTC quantity, $c_t$ the cash balance, and
$C_0 = 10{,}000$ USDT is the initial capital. $\mathrm{ATR}_{168}$ is the
Average True Range \citep{wilder1978atr} computed over a 168-bar (1-week)
window, and $\overline{p}_{t,168}$ is the corresponding simple moving average.
All components are pre-normalized by a 168-bar rolling z-score computed on the
training set and clipped to $[-5, 5]$.

\paragraph{Action $\mathcal{A} = [0,1]^2$.}
At each step the policy emits a continuous 2D action $\mathbf{a}_t =
(a_t^{(0)}, a_t^{(1)})$. The first coordinate $a_t^{(0)}$ (\emph{aggressiveness})
controls two buy limit orders, and $a_t^{(1)}$ (\emph{profit target}) controls
two sell limit orders. We avoid a discrete buy/sell/hold action space because
fine-grained grid width adjustment and volatility adaptation are most naturally
expressed in continuous values; discretization would artificially limit the
expressiveness of the action space.

\paragraph{ATR-proportional grid formula.}
Grid widths and positions are determined through volatility-normalized action
coordinates. With $r_t = \mathrm{ATR}_{168}(p) / p_t$, current price $p_t$, and
average cost $\bar{p}^{\mathrm{avg}}_t$,
\begin{align}
g^{\mathrm{buy\_hi}}_t  &= r_t \cdot (A_b + B_b \cdot a_t^{(0)}) \\
g^{\mathrm{buy\_lo}}_t  &= r_t \cdot (C_b + D_b \cdot a_t^{(0)}) \\
g^{\mathrm{sell\_mkt}}_t &= r_t \cdot (A_s + B_s \cdot a_t^{(1)}) \\
g^{\mathrm{sell\_cost}}_t &= r_t \cdot (C_s + D_s \cdot a_t^{(1)})
\end{align}
Quote prices:
\begin{equation}
q^{\mathrm{buy\_hi}}_t = p_t (1 - g^{\mathrm{buy\_hi}}_t), \quad
q^{\mathrm{buy\_lo}}_t = p_t (1 - g^{\mathrm{buy\_lo}}_t)
\end{equation}
\begin{equation}
q^{\mathrm{sell\_mkt}}_t = p_t (1 + g^{\mathrm{sell\_mkt}}_t), \quad
q^{\mathrm{sell\_cost}}_t = \bar{p}^{\mathrm{avg}}_t (1 + g^{\mathrm{sell\_cost}}_t)
\end{equation}
The coefficients $\{A_b, B_b, C_b, D_b, A_s, B_s, C_s, D_s\}$ are fixed
environment constants, with values reported in \S\ref{sec:atr_baseline}. Only
the \texttt{sell\_cost} quote is anchored at the average cost; the choice
explicitly exposes ``break-even recovery'' to the policy.

\paragraph{Reward.}
The baseline (\texttt{sym}, symmetric) form is the 1-step PnL ratio:
\begin{equation}
r_t^{\mathrm{sym}} = \frac{E_t - E_{t-1}}{C_0}, \quad
E_t = c_t + h_t \cdot p_t
\end{equation}
Transaction fees are already deducted from $c_t$ inside
\texttt{\_execute\_buy} and \texttt{\_execute\_sell}, so no separate subtraction
appears in the reward. All four reward variants in this paper are defined as
functions of $r_t^{\mathrm{sym}}$ and detailed in \S\ref{sec:reward_variants}.

\subsection{Trading-Environment Dynamics}
\label{sec:env_dynamics}

This subsection details order execution, the cycle definition, and the
transaction-cost model.

\paragraph{Limit-order execution.}
Fill decisions for the four quotes use the next bar's high/low
$(p^{\mathrm{H}}_{t+1}, p^{\mathrm{L}}_{t+1})$:
\begin{equation}
\begin{aligned}
&p^{\mathrm{L}}_{t+1} \leq q^{\mathrm{buy\_hi}}_t  \Rightarrow
\text{fill at $q^{\mathrm{buy\_hi}}_t$} \\
&p^{\mathrm{L}}_{t+1} \leq q^{\mathrm{buy\_lo}}_t  \Rightarrow
\text{fill at $q^{\mathrm{buy\_lo}}_t$} \\
&p^{\mathrm{H}}_{t+1} \geq q^{\mathrm{sell\_mkt}}_t \Rightarrow
\text{fill at $q^{\mathrm{sell\_mkt}}_t$} \\
&p^{\mathrm{H}}_{t+1} \geq q^{\mathrm{sell\_cost}}_t \Rightarrow
\text{fill at $q^{\mathrm{sell\_cost}}_t$}
\end{aligned}
\end{equation}
When both buy and sell quotes are filled in the same bar we apply a
\emph{sell-first} rule, prioritizing inventory liquidation. The fill price is
the \emph{quote price itself}, not the next-bar extreme. Filling at the
next-bar extreme injects favorable bias into the simulation; we therefore
discarded an earlier ``next\_low/next\_high'' execution rule and use only the
limit-order rule throughout this paper.

\paragraph{Cycle definition.}
A cycle starts when the \emph{first buy fills while holdings = 0} and ends
when subsequent sells return holdings to 0. Multiple partial buys/sells can
occur inside a single cycle, and the cycle return is
$\mathrm{cycle\_pnl} = (c_{\text{end}} - c_{\text{start}}) / c_{\text{start}}$.
At the start of each cycle, the slot size used for the cycle is set:
\begin{equation}
\text{slot\_size}_{\text{cycle}} =
\frac{c_{\text{start}}}{n_{\text{splits}}}, \quad
\text{order\_size} = \frac{\text{slot\_size}}{n_{\text{buy\_orders}}}
\end{equation}
We use $n_{\text{splits}}=2$ and $n_{\text{buy\_orders}}=2$ (the best values
from Phase 2 Optuna Trial \#34). \texttt{threshold\_basis} is fixed to
\texttt{avg\_price}, meaning the ``liquidate all'' detection uses the average
cost.

\paragraph{Transaction-cost model.}
Transaction costs are set to \emph{0.05\% per side} (Binance maker fee). On a
buy, the filled quantity is $\Delta h = \text{order\_size} \cdot (1 - f) / q$
and the cash change is $\Delta c = -\text{order\_size}$, so the fee is already
reflected in $c_t$. In \S\ref{sec:slippage} (exp033) we add an additional
slippage $\sigma = 0.02\%$ that modifies fill prices to
$q^{\mathrm{buy}} \cdot (1 + \sigma)$ and $q^{\mathrm{sell}} \cdot (1 - \sigma)$.
Slippage is not applied to the base environment and appears only in the
robustness analysis.

\subsection{ATR Baseline}
\label{sec:atr_baseline}

The comparison baseline \emph{shares the environment and grid formula} with
our RL policies but determines actions via a rule. Concretely, the ATR
baseline reuses the grid formulas of \S\ref{sec:mdp} for $g^{\mathrm{buy\_hi}},
g^{\mathrm{buy\_lo}}, g^{\mathrm{sell\_mkt}}, g^{\mathrm{sell\_cost}}$ with all
RL action degrees of freedom \emph{zeroed} ($a_t^{(0)} = a_t^{(1)} = 0$, so
that the $B_b, D_b, B_s, D_s$ terms drop and the grid widths are fixed
time-independent constants determined entirely by $A_b, C_b, A_s, C_s$). The
coefficient values come from Phase 2 Optuna Trial \#34:

\begin{table}[h]
\centering
\caption{Grid coefficients of the ATR baseline (Phase 2, Optuna Trial \#34).}
\begin{tabular}{lcclcc}
\toprule
Coef. & Value & Role & Coef. & Value & Role \\
\midrule
$A_b$ & 1.665 & buy\_hi base width & $A_s$ & 0.285 & sell\_mkt base width \\
$C_b$ & 6.070 & buy\_lo base width & $C_s$ & 1.951 & sell\_cost base width \\
$B_b$ & 1.748 & buy\_hi action sensitivity & $B_s$ & 1.950 & sell\_mkt action sens.\ \\
$D_b$ & 18.683 & buy\_lo action sensitivity & $D_s$ & 7.500 & sell\_cost action sens.\ \\
\bottomrule
\end{tabular}
\label{tab:atr_coefs}
\end{table}

The ATR baseline yields Val Sharpe $1.378$ without slippage,\footnote{In the
early RESEARCH\_LOG (May 14, 2026) we recorded $1.505$; the value was
re-measured to $1.378$ after the Phase 16a unification of the evaluation
setup ($n_{\text{eval\_episodes}}$, env config synchronization). We use
$1.378$ throughout the paper.} dropping to $0.835$ when slippage of $0.02\%$
is added.

\subsection{Four Reward Variants}
\label{sec:reward_variants}

We now define the four reward variants that form the principal comparison unit
of this paper. All are expressed as functions of $x_t := r_t^{\mathrm{sym}} =
(E_t - E_{t-1}) / C_0$ and differ \emph{only} in functional form; the state,
action, ATR grid, and fee model are otherwise identical.

\paragraph{(i) Symmetric (\texttt{sym}).}
\begin{equation}
r_t^{\mathrm{sym}} = x_t = \frac{E_t - E_{t-1}}{C_0}
\end{equation}
No hyperparameter. Baseline of this paper.

\paragraph{(ii) Asymmetric (\texttt{asym}).}
\begin{equation}
r_t^{\mathrm{asym}} =
\begin{cases}
x_t & \text{if } x_t \geq 0 \\
\beta \cdot x_t & \text{if } x_t < 0
\end{cases}, \quad \beta > 1
\end{equation}
A first-order approximation of a risk-averse utility, weighting losses by
$\beta$. An Optuna 30-trial $\times$ 200k-step search over the range
$[1.0, 4.0]$ to maximize Val Sharpe selected the best $\beta = 3.420$.

\paragraph{(iii) Differential Sharpe Ratio (\texttt{dsr}).}
The formulation of \citet{moody2001dsr}. Given an EWMA decay rate $\eta$,
online estimates of the first and second return moments evolve as
\begin{align}
A_t &= A_{t-1} + \eta \, (x_t - A_{t-1}) \\
B_t &= B_{t-1} + \eta \, (x_t^2 - B_{t-1}) \\
\Delta A_t &= x_t - A_{t-1}, \quad \Delta B_t = x_t^2 - B_{t-1}
\end{align}
The DSR step reward is the differential of the on-line Sharpe estimator
$\hat{S}_t = A_t / \sqrt{B_t - A_t^2}$:
\begin{equation}
r_t^{\mathrm{dsr}} =
\frac{B_{t-1} \, \Delta A_t - \tfrac{1}{2} A_{t-1} \, \Delta B_t}
     {(B_{t-1} - A_{t-1}^2)^{3/2}}
\end{equation}
A 30-trial Optuna search over $[1/720, 1/24]$ (log-uniform) selected the best
$\eta = 0.0352 \approx 1/28\,\text{h}$, an EWMA window of approximately
1.2~days.

\paragraph{(iv) Prospect-theoretic (\texttt{pt}).}
The utility function of \citet{kahneman1979prospect} and
\citet{tversky1992cumulative} applied directly as a step reward:
\begin{equation}
r_t^{\mathrm{pt}} =
\begin{cases}
\mathrm{sgn}(x_t) \, |x_t|^{\alpha} & \text{if } x_t \geq 0 \\
-\lambda \, |x_t|^{\alpha} & \text{if } x_t < 0
\end{cases}, \quad 0 < \alpha \leq 1, \quad \lambda \geq 1
\end{equation}
$\alpha < 1$ encodes the concavity of the utility (diminishing marginal
utility of gains), and $\lambda > 1$ encodes loss aversion. The standard
values from \citet{tversky1992cumulative}, fitted to human behavior, are
$\alpha = 0.88$ and $\lambda = 2.25$. The Val-Sharpe-maximizing values for our
RL policy, however, are $\alpha = 0.683$ and $\lambda = 3.303$, found by
Optuna 30 trials over $\alpha \in [0.5, 1.0]$ and $\lambda \in [1.0, 4.0]$.
The fact that RL prefers a \emph{more concave} utility and \emph{stronger
loss aversion} than the human standard is consistent with the mechanism in
\S\ref{sec:mechanism} and the OOS robustness reported in \S\ref{sec:oos}.

\paragraph{Optuna search summary.}
For each of the three variants \texttt{asym}, \texttt{dsr}, \texttt{pt} we ran
30 trials $\times$ 200k steps using a TPE sampler
\citep{akiba2019optuna} (seed=42, MedianPruner with
$n_{\text{startup}}=5$). The total training cost was $3 \times 30 \times
200\text{k} = 18$M steps, taking approximately 2 hours on a Windows 11
CPU with $n_{\text{envs}}=4$. \texttt{sym} has no hyperparameter and skips
this stage.

\subsection{PPO Training and the Stabilization Package}
\label{sec:ppo_stab}

Policy training uses PPO \citep{schulman2017ppo} from stable-baselines3
\citep{raffin2021sb3}. Standard PPO hyperparameters
($\gamma=0.99$, $\lambda_{\mathrm{GAE}}=0.95$, clip
$\epsilon=0.2$, $v_f$ coef $=0.5$, gradient clip $=0.5$, MLP policy
$[64, 64]$, batch=256, $n_{\text{epochs}}=10$, $n_{\text{steps}}=4096$,
$n_{\text{envs}}=4$) follow stable-baselines3 defaults. To mitigate the
reward sparsity of our environment (the per-step reward is near 0 when no
fill occurs) and the late-stage policy collapse observed in Phase 2
(exp020--022), we add the following \emph{stabilization package} (exp030):

\begin{enumerate}[leftmargin=*]
\item \textbf{Learning-rate linear decay}: $\alpha_t$ linearly decays from
$3 \times 10^{-4}$ to $1 \times 10^{-5}$ over training. This bounds the size
of late-stage policy updates and dampens post-best collapse.
\item \textbf{Entropy coefficient annealing}: $c_{\text{ent}}$ linearly anneals
from $0.01$ to $0.001$, ensuring sufficient initial exploration while
encouraging convergence to a near-deterministic policy late in training.
\item \textbf{Target KL early stop}: $\text{target\_kl}=0.02$ measures the KL
divergence after each minibatch update and aborts the update for that batch if
the threshold is exceeded, preventing PPO from making myopic large policy
changes.
\item \textbf{Best-model checkpoint}: Val Sharpe is evaluated every 50k steps
($n_{\text{eval\_episodes}}=5$) and the weights at the best Val Sharpe are
saved separately, providing protection against late-stage collapse.
\end{enumerate}

The package was developed in exp030 (with \texttt{sym} reward) and produced
Val Sharpe best $1.974$ at step 550k, falling to final $1.209$ at step 1M.
Because the collapse pattern after step 700k persists despite the package,
\emph{all evaluations in this paper use the best checkpoint}; the final model
is never used.

\subsection{Evaluation Protocol}
\label{sec:eval_protocol}

We use \emph{three evaluation environments concurrently}. The data partition
is strict time-order to prevent leakage:

\begin{table}[h]
\centering
\caption{Data partition.}
\begin{tabular}{lll}
\toprule
Partition & Range & Bars (approx.) \\
\midrule
Train & 2017-08-17 \,--\, 2020-12-31 & $\sim$30{,}000 \\
Val (single-split) & 2021-01-01 \,--\, 2023-12-31 & $\sim$26{,}000 \\
Test (sealed) & 2024-01-01 \,--\, 2026-04 & $\sim$20{,}000 \\
\bottomrule
\end{tabular}
\label{tab:data_split}
\end{table}

\paragraph{Val (single-split).}
The main comparison environment of \S\ref{sec:pareto}. For each reward variant
we train 10 seeds $\times$ 1M steps and evaluate the best\_model on the entire
Val range with $n_{\text{eval\_episodes}}=5$, producing per-seed
Sharpe/MDD/Calmar/trade-count/return.

\paragraph{CPCV 6-fold.}
The multi-split environment of \S\ref{sec:cpcv}. Following
\citet{lopezdeprado2018advances}, we apply 6-fold Combinatorial Purged
Cross-Validation \emph{inside the Train partition}, producing
$\binom{6}{2}=15$ paths (each consisting of 5 months of training plus
1 month of evaluation). A 1-week purge gap between train and evaluation folds
prevents information leakage. We execute 4 variants $\times$ 15 paths $\times$
1 seed = 60 runs.

\paragraph{Test (unsealed).}
The OOS environment of \S\ref{sec:oos}. For honesty, following
\citet{gort2022backtest}, Test was \emph{kept sealed from any search or
hyperparameter decision} until exp035 (May 16, 2026). After unsealing, we
evaluate 100 RL models (4 variants $\times$ multi-seed/seed-pair) along with
the ATR baseline and Buy-and-Hold on Test.

\paragraph{Statistical tests.}
All pairwise comparisons report Cohen's $d$ and bootstrap $\mathrm{P}(A > B)$
($10^4$ resamples). Single-sided t-tests against ATR on CPCV use Bonferroni
4-way correction along with the Deflated Sharpe Ratio test of
\citet{lopezdeprado2014dsr}. To report evaluation variability and distributional
skew, we follow \citet{henderson2018rlreproducibility} and add IQM
(Interquartile Mean) and 5\% CVaR as secondary metrics.

% ─── 4. Phase 1 Negative Finding ─────────────────────────────────────────────
\section{Phase 1: Policy Saturation under Symmetric Reward}
\label{sec:phase1_negative}

\textit{[TBD. Reproduces the ``RL $=$ Fixed $[1.0, 0.0]$'' finding from
exp020--022 of the earlier phase. Motivates reward formulation in this paper.]}

% ─── 5. Main Finding: Pareto Frontier ────────────────────────────────────────
\section{Pareto Frontier and Scenario D}
\label{sec:pareto}

\textit{[TBD. exp032b results:
4 variants $\times$ 10 seeds, pairwise statistics, Pareto scatter,
two-cluster discovery, hypothesis H1--H3 evaluation,
Scenario D confirmation.]}

% ─── 6. Mechanism Quantification ─────────────────────────────────────────────
\section{Mechanism Quantification}
\label{sec:mechanism}

\textit{[TBD. exp032c:
policy distance matrix, behavior per regime, counterfactual action map,
action distribution; causal chain reward $\to$ trade frequency $\to$ cluster.]}

% ─── 7. Robustness ───────────────────────────────────────────────────────────
\section{Robustness Analysis}
\label{sec:robustness}

\subsection{Slippage Robustness}
\label{sec:slippage}

\textit{[TBD. exp033 results with slippage 0.02\%.
Cluster preservation ratio 2.19$\times$;
ATR-with-slippage fair-comparison results (Phase 16a).]}

\subsection{Multi-Split Evaluation via CPCV}
\label{sec:cpcv}

\textit{[TBD. exp034: 6-fold CPCV 15 paths.
DSR $p<0.001$ (Bonferroni 4-way correction);
\texttt{dsr} reversal to 1\textsuperscript{st}
(1.413, 5\% CVaR 0.890); cluster preservation ratio 3.55$\times$.]}

\subsection{Out-of-Sample Evaluation and Prospect Theory's Robustness}
\label{sec:oos}

\textit{[TBD. Central chapter.
exp035: 100 RL models + ATR baseline + Buy-and-Hold on Test.
\texttt{pt} 1\textsuperscript{st} on Test (both sources $p<0.002$);
\texttt{dsr} last on Test.
Phase 16d mechanism: hold-duration distribution
(DSR 7-day max vs.\ pt 6h max);
causal account of \texttt{pt}'s OOS robustness.]}

% ─── 8. Discussion ───────────────────────────────────────────────────────────
\section{Discussion}
\label{sec:discussion}

\textit{[TBD. Subsections:
\S8.1 Val--Test distribution-shift quantification (KS $p<10^{-10}$, $-27\%$
variance);
\S8.2 honest acknowledgment of environment dependence in exp027\_rl prior
evidence (Env-v3 \texttt{asym} 1.955);
\S8.3 limits of single-metric winner;
\S8.4 limitations: BTC-only, single regime in Test, single slippage level.]}

% ─── 9. Conclusion ───────────────────────────────────────────────────────────
\section{Conclusion}
\label{sec:conclusion}

\textit{[TBD. Four main messages:
(1) Pareto-frontier trade-off;
(2) winner reversal across evaluation environments;
(3) OOS robustness of Prospect Theory;
(4) in-sample advantage $\neq$ OOS robustness.
Future work: multi-asset, additional regimes, DR, meta-RL.]}

% ─── References ──────────────────────────────────────────────────────────────
\bibliographystyle{plainnat}
\bibliography{references}

\end{document}
